\documentclass[12pt]{article}

\usepackage{problem_set_style}
\usepackage{kotex}
\usepackage{booktabs}
\usepackage{array}
\usepackage{graphicx}

\begin{document}
\heading{8}{2027. X. XX (예정)}{}{Fall 2027}{중간 프로젝트 공고 (샘플)}

\section*{1. 프로젝트 개요}
4$\sim$7주차에서 배운 표 기반(tabular) 강화학습 기법(동적계획법,
몬테카를로, 시간차 학습, n-step 부트스트래핑)을 실제 Gymnasium
환경에 적용하는 학기 전반 8주에 걸쳐 진행하는 팀 프로젝트이다. 목표는 ``좋은 정책을
만드는 것'' 자체가 아니라, \textbf{그 성능 숫자를 정직하게 얻는
절차}를 실제로 수행하고 보고하는 것이다.

\vspace{4pt}
\noindent 진행 순서: \textbf{환경 고르기 $\to$ Abstract $\to$ 구현 $\to$ 평가 $\to$ 보고서 제출 $\to$ 동료평가}.

\section*{2. 필수 요구사항}
\begin{itemize}
\item \textbf{환경}: Gymnasium의 표 기반(이산 상태·행동) 환경
  하나 --- \texttt{FrozenLake-v1}, \texttt{CliffWalking-v1},
  \texttt{Taxi-v4} 등.
\item \textbf{적용 알고리즘}: 4$\sim$7주차 중 \textbf{최소 두
  가지}(예: Q-learning vs SARSA, MC 제어 vs n-step TD)를 같은
  환경에 적용해 학습 곡선과 최종 정책을 비교한다.
\item \textbf{수식적 정당화 최소 1개}: 예 --- Q-learning/SARSA
  비교라면 on/off-policy 관점에서 정책이 왜 다르게 학습됐는지,
  $\varepsilon$/$\alpha$ 변경 실험이라면 수렴 조건과 연결, Dyna-Q
  라면 계획 스텝 수와 학습 속도의 관계를 직접 측정.
\item \textbf{베이스라인 필수}: 무작위 정책의 평균 리턴을 먼저
  측정해, ``배운 정책이 실제로 무작위보다 나은가''를 판정할
  기준점을 마련한다.
\item \textbf{시드 $\ge 3$}, 모든 시드에 \textbf{동일한 평가
  프로토콜}(greedy 평가 $\varepsilon=0$, 동일 에피소드 수)을
  적용한다.
\end{itemize}

\newpage
\section*{3. 환경 후보 (예시)}
\begin{center}
\small
\begin{tabular}{@{}lllp{4.4cm}@{}}
\toprule
환경 & 상태$\times$행동 & 보상 구조 & 잘 어울리는 비교 \\
\midrule
FrozenLake-v1 & 16$\times$4 (미끄럼) & 도착 +1, 그 외 0 (구멍 종료)
& MC vs SARSA --- 확률 전이에서 시드 안정성 \\
CliffWalking-v1 & 48$\times$4 & 스텝 $-1$, 절벽 $-100$, 도착 0 &
SARSA vs Q-learning --- on/off-policy 클래식 비교 \\
Taxi-v4 & 500$\times$6 & 성공 +20, 불법 $-10$, 스텝 $-1$ &
n-step vs TD(0), Dyna-Q \\
\bottomrule
\end{tabular}
\end{center}
\vspace{4pt}
환경 선택은 자유이지만, \textbf{4$\sim$7주차의 개념 검증
목적}에 근거해야 한다 (상태 수, 보상 구조, 확률 전이 중 무엇이 이
비교를 정당화하는지 설명할 수 있어야 한다).

\section*{4. 시드 프로토콜}
``한 번의 실행''은 아무것도 말해주지 않는다 --- $\varepsilon$-greedy
탐험의 무작위성, 환경 전이 자체의 무작위성, 정책 초기화 때문에
같은 코드·같은 파라미터를 그대로 돌려도 매번 다른 학습 곡선이
나온다.
\begin{enumerate}
\item \textbf{시드 $\ge 3$}으로 비교할 모든 알고리즘을 전부
  돌린다.
\item 학습 도중 성과(탐험 포함)와 학습 후 \textbf{greedy 평가}
  성과를 반드시 분리해서 보고한다.
\item 하이퍼파라미터는 학습 곡선(validation)으로만 튜닝하고,
  최종 평가 수치를 보며 조정하지 않는다.
\end{enumerate}

\newpage
\section*{5. 일정 (전반 8주)}
\begin{center}
\small
\begin{tabular}{@{}p{3.6cm}p{5.4cm}p{3.7cm}@{}}
\toprule
\textbf{시기} & \textbf{해야 할 일} & \textbf{제출} \\
\midrule
\textbf{1$\sim$2주} & 팀 구성, 환경 후보 2$\sim$3개 조사(상태 수·보상 구조·확률 전이 여부) &  \\
\textbf{3주} & 문제 정의(MDP 5요소, 주 지표), 선택한 환경, 무작위 정책 베이스라인과 비교할 알고리즘 계획 & \textbf{Abstract} (1p) \\
\textbf{4$\sim$5주} & 두 알고리즘 구현, 시드 3개 이상 실행, $\varepsilon$/$\alpha$ 튜닝(모두 동일 프로토콜), greedy 평가, 시드 비교표 &  \\
\textbf{6주 (중간고사 기간)} & 보고서 완성 & \textbf{보고서 제출}: 보고서 PDF + 코드 + 학습 곡선 \\
\textbf{7주} & 제출된 다른 팀 보고서 읽기 &  \\
\textbf{8주 (중간고사 기간)} & 다른 팀 2개에 대한 peer-review 작성 & \textbf{동료평가} 제출 (리뷰 2건) \\
\bottomrule
\end{tabular}
\end{center}
\vspace{4pt}
Abstract에 무작위 정책 베이스라인 계획을 넣는 이유: 이 숫자를 봐야 ``배운 정책이 무작위보다 나은가''를 판정할 수 있다.

\section*{6. 보고서 구조}
\begin{enumerate}
\item \textbf{문제 정의} (0.5p): MDP 5요소(상태·행동·전이·보상·
  할인율), 주 지표(greedy 평가 평균 리턴), 베이스라인(무작위
  정책 리턴).
\item \textbf{알고리즘 선택과 이유} (1p): 후보 2개 이상, 왜 이
  알고리즘을 이 환경에 골랐는가(보상 구조·상태 수와 연결),
  하이퍼파라미터($\alpha,\gamma,\varepsilon$ 스케줄)와 선택 근거.
\item \textbf{실험 프로토콜} (1p): 시드 목록, 학습 에피소드 수,
  평가 프로토콜($\varepsilon=0$, 에피소드 수, 모든 알고리즘 동일)
  명시.
\item \textbf{결과} (1p): 시드 비교표, 학습 곡선 1장(여러 시드를
  한 그림에 얇은 선으로), greedy 정책 경로 시각화.
\item \textbf{수식적 정당화} (1p): 배운 개념 1개 이상으로 ``왜
  이 결과''.
\item \textbf{재현 방법} (0.5p): 환경 버전, 시드 목록, 에피소드 수, 하이퍼파라미터, 실행 명령. 다른 사람이 같은 곡선을 다시 얻을 수 있어야 한다.
\item \textbf{한계와 잘 안 된 부분} (0.5p, \textbf{필수}):
  안 된 부분이 없었다면 무엇을 시도했다가 안 됐는지를 쓴다.
\end{enumerate}

\newpage
\section*{7. 채점 기준 (Rubric)}
다섯 항목을 각 20\%로 채점한다. 목적은 ``좋은 정책''이 아니라 \textbf{주장을 정직하고 재현 가능하게 입증}하는 것이다.
\begin{center}
\small
\begin{tabular}{@{}p{3.3cm}cp{7.4cm}@{}}
\toprule
항목 & 비중 & 좋은 답 \\
\midrule
\textbf{1. 주장의 명확성} (clarity) & 20\% & 무엇을 주장하는지 첫 문단에서 한 문장으로 알 수 있고, MDP 5요소·주 지표가 명확하며, 결과가 주장과 이어짐 \\
\textbf{2. 비교 대상} (baseline) & 20\% & 무작위 정책의 평균 리턴과 알고리즘 두 가지 이상이 같은 시드·같은 평가 프로토콜로 비교됨. 비교 없는 숫자는 해석하지 않음 \\
\textbf{3. 확장성} (scalability) & 20\% & 다른 시드 묶음이나 다른 환경(예: 다른 맵 크기) 하나에서도 결론이 유지되는지 확인하고, 안 통하는 조건을 적음 \\
\textbf{4. 재현성} (reproducibility) & 20\% & 환경 버전·시드 목록·에피소드 수·평가 프로토콜·하이퍼파라미터·실행 방법이 적혀 있어 다른 사람이 같은 곡선을 다시 얻을 수 있음 \\
\textbf{5. 한계의 솔직함} (limitations) & 20\% & 불안정했던 시드·구간, 약한 가정, 통하지 않을 환경 조건을 숨기지 않고 원인을 진단함. 성능이 낮아도 한계가 정직하면 높은 점수 \\
\bottomrule
\end{tabular}
\end{center}
\vspace{4pt}
시드 하나만 보여 주거나, 무작위 정책 비교가 없거나, 재현이 불가능하면 크게 감점받는다.

\section*{8. 동료평가 (Peer-Review)}
\begin{itemize}
\item \textbf{Peer-Review}: 8주차에 다른 두 팀의 보고서에 대해 작성 리뷰를 각 1통씩 제출한다. 각 리뷰는 아래 5항목 체크리스트와, 반박 가능한 질문을 1개 이상 포함해야 한다.
\end{itemize}

\vspace{4pt}
\noindent\textbf{Peer-Review 체크리스트 (채점 기준과 같은 5항목)}
\begin{center}
\small
\begin{tabular}{@{}ll@{}}
\toprule
항목 & 확인할 것 \\
\midrule
주장의 명확성 (clarity) & 무엇을 주장하는지 분명하게 썼는가 \\
비교 대상 (baseline) & 비교 대상이 있는가 (무작위 정책 포함) \\
확장성 (scalability) & 다른 환경에도 적용 가능한가 \\
재현성 (reproducibility) & 다른 사람이 다시 해 볼 수 있는가 \\
한계 (limitations) & 한계를 솔직하게 적었는가 \\
\bottomrule
\end{tabular}
\end{center}

\section*{9. 팀 구성 및 제출 방식}
\begin{itemize}
\item \textbf{팀 구성}: 팀당 인원 및 팀 편성 방식은 추후 공지.
\item \textbf{제출물}: Abstract(3주), 보고서(PDF)와 코드(6주), peer-review 2건(8주).
\item \textbf{제출 방법 및 마감}: 추후 공지 (LMS/과제 제출 시스템
  및 마감 일시는 별도 안내 예정).
\end{itemize}

\end{document}
